%! AP Statistics — Unit 8, Lesson 8.3 — Group Activity (Answer Key)
\documentclass[10pt]{article}
\usepackage{apstats-article}
\usepackage{apstats-key}

\begin{document}

\pageheader{Unit 8, Lesson 8.3}{Group Activity --- Answer Key}
\namepartnerperiod

\vspace{0.05in}

\begin{tcolorbox}[colback=sky!30, colframe=navy, boxrule=0.8pt, arc=2mm,
    left=3mm, right=3mm, top=2mm, bottom=2mm]
\textbf{Birth Month Distribution of MLB Players}\\[4pt]
Researchers wonder whether Major League Baseball (MLB) players are more likely to be born
in certain months. If birth month had no effect on player development, we would expect
players to be born equally often in each month ($p = 1/12 \approx 0.0833$ for each month).
A random sample of 180 MLB players yields the following birth month counts:

\renewcommand{\arraystretch}{1.4}
\begin{center}
\begin{tabular}{lc|lc}
    \textbf{Month} & \textbf{Observed} & \textbf{Month} & \textbf{Observed} \\
    \hline
    January   & 21 & July      & 13 \\
    February  & 18 & August    & 12 \\
    March     & 19 & September & 16 \\
    April     & 22 & October   & 14 \\
    May       & 17 & November  & 11 \\
    June      & 15 & December  & 22 \\
\end{tabular}
\end{center}
\end{tcolorbox}

\vspace{0.08in}

\begin{enumerate}[leftmargin=1.5em]

  \item \textbf{Hypotheses.}

        $H_0$: \ans{The birth month distribution of MLB players is uniform; each month
        accounts for $1/12$ of all players ($p_1 = p_2 = \cdots = p_{12} = 1/12$).}

        \vspace{0.05in}
        $H_a$: \ans{At least one month has a different proportion of MLB player birthdays
        than $1/12$.}

        \vspace{0.08in}

  \item \textbf{Expected Counts.}

        $E = $ \ans{$180 \times (1/12) = 15$} for each month.

        \vspace{0.05in}
        Large-counts condition: \ans{All expected counts $= 15 \ge 5$; condition is satisfied.}

        \vspace{0.08in}

  \item \textbf{Test Statistic.}

\renewcommand{\arraystretch}{2.0}
\begin{center}
\begin{tabular}{lcccc}
    \textbf{Month} & $O$ & $E$ & $(O-E)^2/E$ & \\
    \hline
    January   & 21 & \ans{15} & \ans{$36/15 = 2.400$} & \\
    February  & 18 & \ans{15} & \ans{$9/15 = 0.600$}  & \\
    March     & 19 & \ans{15} & \ans{$16/15 \approx 1.067$} & \\
    April     & 22 & \ans{15} & \ans{$49/15 \approx 3.267$} & \\
    May       & 17 & \ans{15} & \ans{$4/15 \approx 0.267$}  & \\
    June      & 15 & \ans{15} & \ans{$0/15 = 0.000$}  & \\
    July      & 13 & \ans{15} & \ans{$4/15 \approx 0.267$}  & \\
    August    & 12 & \ans{15} & \ans{$9/15 = 0.600$}  & \\
    September & 16 & \ans{15} & \ans{$1/15 \approx 0.067$}  & \\
    October   & 14 & \ans{15} & \ans{$1/15 \approx 0.067$}  & \\
    November  & 11 & \ans{15} & \ans{$16/15 \approx 1.067$} & \\
    December  & 22 & \ans{15} & \ans{$49/15 \approx 3.267$} & \\
    \hline
    \textbf{Total} & 180 & & \ans{$\approx 12.93$} $= \chi^2$ & \\
\end{tabular}
\end{center}

  \item \textbf{Degrees of Freedom and P-Value.}

        $df =$ \ans{$12 - 1 = 11$}

        \vspace{0.05in}
        \ansline{The probability, given that birth month is uniformly distributed among all MLB players, of obtaining a chi-square statistic of approximately 12.93 or larger is about 0.30.}

  \item \textbf{Conclusion.}

        \ansline{Because $p \approx 0.30 > \alpha = 0.05$, we fail to reject $H_0$. We do not have convincing evidence that the birth month distribution of MLB players differs from a uniform distribution.}

\end{enumerate}

\end{document}
